\documentclass[10pt,a4paper]{amsart}
\usepackage[margin=19mm]{geometry}
\usepackage{amsmath,amssymb,mathtools}
\usepackage[scr=rsfs,cal=euler]{mathalfa}
\usepackage{xfrac}
\usepackage[unicode,hidelinks,bookmarks=false]{hyperref}
\newcommand{\ie}{i.e.\ }
\newcommand{\cf}{cf.\ }
\newcommand{\resp}{respectively\ }
\newcommand{\Subsection}{\subsection}
\providecommand{\nobreakdash}{\nobreak\hbox{-}}
\setlength{\emergencystretch}{2em}
\multlinegap0pt
\usepackage{mdframed}
\usepackage{mdframed}
\usepackage{mdframed}
\usepackage{mdframed}
\usepackage{amssymb}
\usepackage{mdframed}
\newtheorem{theorem}{Theorem}
\newtheorem{claim}{Claim}
\renewcommand{\P}{\bb{P}}
\newcommand{\bb}{\mathbb}
\newcommand{\dom}{\mathrm{dom}}
\title{Preservation of some topological properties under forcing}
\author{Chris Lambie-Hanson, Pedro Marun}
\date{Source: arXiv:2601.07624v1 (CC BY 4.0)}
\begin{document}
\maketitle

\noindent\emph{Source and licence.} Preservation of some topological properties under forcing. Version arXiv:2601.07624v1 (CC BY 4.0), distributed by arXiv under the Creative Commons Attribution 4.0 International licence, http://creativecommons.org/licenses/by/4.0/, as recorded in the arXiv licence metadata for this exact version and shown on the abstract page https://arxiv.org/abs/2601.07624v1. Full licence notice: \texttt{licenses/arxiv-2601.07624-CC-BY-4.0.txt}.\par\smallskip

\noindent\textit{Editorial scope.} Counted unit: the article's theorem main (source0.tex lines 381--389). The article's complete original proof of it is delivered with it (source0.tex lines 391--495, one \texttt{proof} environment of 105 lines). No mathematics is written, repaired or re-derived: the statement and the proof are the article's own lines, in the article's own notation, and no retained line is substituted or re-flowed.  No further source-local context is needed: every cross-reference of every retained line resolves inside the two delivered spans.  Nothing was omitted from any retained span, so the package carries no omission record.  No retained line contains a citation, so the wrapper carries no bibliography and no bibliography entry.  No retained line contains a figure environment, an image-inclusion command or drawing code, so no image is delivered and the package declares no asset dependency.  Editorial formatting only, in addition: an AMS \texttt{amsart} preamble that restates the article's own declarations and macros used by the retained lines, namely the environments \texttt{theorem}, \texttt{claim} on separate counters, together with the standard packages those lines need.  The retained environments print in this document as Theorem 1, Claim 1 in order of appearance, whereas the article prints the same items with the numbering of its own full document.  The complete native source of the article as distributed in the arXiv e-print for this exact version is bundled unchanged as \texttt{sources/192518/source0.tex}.
\begin{theorem}\label{main}
  Suppose that $(X, \tau)$ is a topological space and player II has a winning strategy in one of the following games:
\begin{itemize}
\item the strong countable fan tightness game for $X$ at $x\in X$;
\item the Rothberger game for $X$;
%\item the Menger game for $X$.
\end{itemize}
Then this is preserved in any forcing extension of $V$.
\end{theorem}

\begin{proof}
We will prove the theorem for the strong countable fan tightness game, the proof for the Rothberger game is similar and left to the reader.

  Let $\P$ be an arbitrary forcing notion. Given a point $x \in X$, let $\Game_x(X)$ denote the 
  strong countable fan tightness game for $X$ at the point $x$. Recall that $\Game_x(X)$ is a two-player game of length $\omega$. In each round, player I first plays a set $A_n \subseteq X$ such that $x \in \mathrm{cl}(A_n)$; then player II plays a point $x_n \in A_n$. Player II wins if $x \in \mathrm{cl}(\{x_n :n <\omega\})$; otherwise, player I wins.
  
  Let $\sigma_x$ be a winning strategy for player II in $\Game_x(X)$. Concretely, 
  $\sigma_x$ is a function such that $\dom(\sigma_x)$ is the set of all finite sequences 
  $\vec{A} = \langle A_k :k \leq n \rangle$ such that each $A_k$ is a subset of $X$ with 
  $x \in \mathrm{cl}(A_k)$. For each such $\vec{A}$, we have $\sigma(\vec{A}) \in A_n$.
  
  We aim to show that player II continues to have a winning strategy in $\Game_x(X)$ in 
  $V^{\P}$. Suppose that $G \subseteq \P$ is generic over $V$, and move to $V[G]$. We will 
  define a winning strategy $\sigma^*$ for player II in $\Game_x(X)$ by recursion on the 
  length of the play. We will ensure that, for a sequence $\vec{A} = \langle A_k :
  k \leq n \rangle$ such that each $A_k$ is a subset of $X$ with $x \in \mathrm{cl}(A_k)$, 
  the strategy outputs not only player II's play $\sigma^*(\vec{A})$ but also auxiliary 
  objects of the following type:
  \begin{itemize}
    \item a $\P$-name $\dot{A}_n \in V$ such that $(\dot{A}_n)_G = A_n$;
    \item a condition $p_n \in G$;
    \item $A^*_n := \{y \in X :\exists p' \leq p_n \ p' \Vdash_{\P} y \in \dot{A}_n\}$.
  \end{itemize}
  To help with the selection of these objects, fix in $V$ a sufficiently large regular cardinal 
  $\theta$ and a wellordering $\vartriangleleft$ of $H(\theta)$.
  
  Suppose that we are given a finite sequence $\vec{A} = \langle A_k :k \leq n \rangle$ 
  of subsets of $X$ such that $x \in \mathrm{cl}(A_k)$ for all $k \leq n$, and suppose that 
  we have already specified $\sigma^*(\vec{A} \restriction m)$ for all $m < n$, as well 
  as auxiliary objects $\dot{A}_k$, $p_k$, and $A^*_k$ for all $k < n$. Assume that we have 
  arranged so that $x \in \mathrm{cl}(A^*_k)$ for all $k < n$.
  
  First, let $\dot{A}_n$ be the $\vartriangleleft$-least $\P$-name for a subset of $X$ such that 
  $(\dot{A}_n)_G = A_n$. Let $q_n$ be the $\vartriangleleft$-least element of $G$ such that, in 
  $V$, we have $q_n \Vdash_{\P} x \in \mathrm{cl}(\dot{A}_n)$. For each $p \leq q_n$, 
  let $A^*_{n,p}$ be the set of all $y \in X$ for which there exists $p' \leq p$ such 
  that $p' \Vdash_{\P} y \in \dot{A}_n$. Note that $A^*_{n,p} \in V$ and 
  $x \in \mathrm{cl}(A^*_{n,p})$.
  
\begin{claim}\label{4.2}
There exists $p \leq q_n$ in $G$ such that 
  $\sigma_x(\langle A^*_k :k < n \rangle ^\frown \langle A^*_{n,p}\rangle) \in A_n$.
\end{claim} 

\begin{proof}
Suppose otherwise. We can then fix a condition $p^*\in G$ such that $p^*\le q_n$ and
\[
p^*\Vdash \neg\exists p\le q_n \left(p\in\dot G\wedge \sigma_x(\langle A^*_k :k < n \rangle ^\frown \langle A^*_{n,p}\rangle) \in \dot{A}_n\right).
\]
Set $y:=\sigma_x(\langle A^*_k :k < n \rangle ^\frown \langle A^*_{n,p^*}\rangle)$. Since $y\in A^*_{n,p^*}$, we can find $p^{**}\le p^*$ such that $p^{**}\Vdash y\in\dot A_n$.

Let $G^*$ be $\P$-generic over $V$ with $p^{**}\in G^*$. Since $p^{**}\le p^*$, it follows that $p^*\in G^*$, hence
\[
V[G^*]\models \neg\exists p\le q_n \left(p\in G^*\wedge \sigma_x(\langle A_k^*:k<n\rangle^\frown \langle A_{n,p}^*\rangle)\in (\dot{A}_n)_{G^*}\right).
\]
On the other hand, $p^{**}\Vdash y\in \dot A_n$, hence
\[
V[G^*]\models \left(p^*\le q_n \wedge p^*\in G^*\wedge \sigma_x(\langle A_k^*:k<n\rangle^\frown \langle A_{n,p^*}^*\rangle)\in (\dot{A}_n)_{G^*}\right).
\]
This is a contradiction.
\end{proof}

Let $p_n$ be the $\vartriangleleft$-least $p$ satisfying the claim, and let 
  $A^*_n := \{y \in X :\exists p' \leq p_n  (p' \Vdash_{\P} y \in \dot{A}_n)\}$. 
  Finally, set $\sigma^*(\vec{A}) = \sigma(\langle A^*_k :k \leq n \rangle)$.
  
  In $V$, let $\dot{\sigma}^*$ be a name for the strategy defined above. 
  We claim that $\dot{\sigma}^*$ is forced to be a winning strategy for player II. 
  Suppose not, and fix in $V$ a sequence $\langle \dot{B}_n :n < \omega \rangle$ 
  and a condition $r \in \P$ such that it is forced by $r$ that
  \begin{itemize}
    \item for all $n < \omega$, $\dot{B}_n$ is a subset of $X$ and $x \in 
    \mathrm{cl}(\dot{B}_n)$;
    \item if $\langle \dot{x}_n :n < \omega \rangle$ is the sequence produced by player II 
    playing $\dot{\sigma}^*$ against $\langle \dot{B}_n :n < \omega \rangle$, 
    then $x \notin \mathrm{cl}(\{\dot{x}_n :n < \omega\})$.
  \end{itemize}
  By extending $r$ further if necessary, we can in fact assume that there is an open set 
  $U \ni x$ such that $r \Vdash_{\P} \{\dot{x}_n :n < \omega\} \cap U = \emptyset$.
  
  For $n < \omega$, let $\dot{a}_n$, $\dot{p}_n$, and $\dot{A}^*_n$ 
  be $\P$-names forced by $r$ to be equal 
  to the auxiliary objects\footnote{We caution the reader that $\dot a_n$ is a name for a name.} $\dot{A}_n$, $p_n$, and $A^*_n$ produced during the run of 
  $\Game_x(X)$ in which player I plays $\langle \dot{B}_n :n < \omega \rangle$ and 
  player II plays according to $\dot{\sigma}^*$. Recursively define a decreasing sequence 
  $\langle r_n :n < \omega \rangle$ of conditions in $\P$ such that, for all $n < \omega$, 
  we have
  \begin{itemize}
    \item $r_n \leq r$;
    \item $r_n$ decides the values of $\dot{a}_n$, $\dot{p}_n$, and $\dot{A}^*_n$, say as 
    $\dot{A}_n$, $p_n$, and $A^*_n$;
    \item $r_n \leq p_n$.
  \end{itemize}
  For each $n < \omega$, note that $A^*_n = \{y \in X :\exists p' \leq p_n \ p' \Vdash_{\P} 
  y \in \dot{A}_n \}$. Now let $\langle x_n :n < \omega \rangle$ be the plays made by player II 
  in a run of $\Game_x(X)$ in which player I plays $\langle A^*_n :n < \omega \rangle$ and 
  player II plays according to $\sigma_x$. Because $\sigma_x$ is a winning strategy for player II, 
  we have $x \in \mathrm{cl}(\{x_n :n < \omega\})$; in particular, there is $n < \omega$ 
  such that $x_n \in U$. Now note that, by our definition of $\dot{\sigma}^*$, $r_n$ forces that 
  $\langle x_k :k \leq n \rangle$ is the sequence of the first $(n+1)$-many moves made 
  by player II in the run of $\Game_x(X)$ in which player I plays $\langle \dot{B}_n :n < 
  \omega \rangle$ and player II plays $\dot{\sigma}^*$; i.e., for all $k \leq n$, we have 
  $r_n \Vdash_{\P} \dot{x}_k = x_k$. Then $x_n \in U$ contradicts the fact that 
  $r_n \leq r$ and $r \Vdash_{\P} \{\dot{x}_i :i < \omega\} \cap U = \emptyset$.
\end{proof}

\end{document}
